\section{Holomorphic--Topological Locality}
\label{sec:HT-operad}

A 3d HT theory places observables in two environments at once:
holomorphic collisions in $\FM_k(\C)$, and ordered topology on
$\Conf_k^<(\R)$. Either alone has a standard local-to-global
principle (factorization algebras on $\C$; $E_1$-algebras on $\R$),
but the coupling is not automatic. A bulk OPE pole at
$z_1 = z_2$ occurs while the boundary ordering on~$\R$ is held
fixed; a rearrangement of the ordering on~$\R$ is performed while
holomorphic positions stay generic. What local-to-global
principle couples these two structures without collapsing either,
and in which operad does the coupling live?

The two-colored Swiss-cheese operad $\SCchtop$ supplies the
answer: its closed colour, the chiral operad $\cP^{\mathrm{ch}}$
built below from logarithmic forms on the Fulton--MacPherson
compactifications $\FM_k(\C)$, governs holomorphic collisions; its
open colour $\Eone$, the little-intervals operad, governs
topological ordering; and the Swiss-cheese composition constraint
(no open-to-closed map) enforces bulk-to-boundary directionality.
An $\Ainf$~chiral algebra is an algebra over
$C_\ast(W(\SCchtop))$; the boundary algebra $A_b =
\End_\cC(b)$ of the open-sector factorization category is the
universal example.

\subsection{The Chiral Operad}
\label{sec:chiral_operad}

\subsubsection{Configuration Spaces and Fulton--MacPherson Compactifications}

A chiral operation on $n$ insertions is controlled by the relative
positions and collision rates of those insertions on~$\C$. The
space $\Conf_n(\C)$ records the positions but is noncompact, so its
top-degree forms do not integrate; the Fulton--MacPherson
compactification $\FM_n(\C)$ adds boundary strata encoding the
relative rates at which points approach each other, making the
logarithmic-residue calculus that extracts OPE coefficients
well-defined.

\begin{definition}[Configuration Spaces]
For $n \geq 1$, the \emph{configuration space} of $n$ distinct ordered points in $\C$ is
\[
\Conf_n(\C) := \{(z_1,\ldots,z_n) \in \C^n \mid z_i \neq z_j \text{ for } i \neq j\},
\]
equipped with the subspace topology from $\C^n \cong \R^{2n}$.
\end{definition}

The space $\Conf_n(\C)$ is a smooth noncompact manifold of real dimension $2n$. Extracting chiral OPE residues by integration requires a compactification that records the relative rates at which points approach each other. The Fulton--MacPherson compactification \cite{FM94} is the canonical choice:

\begin{definition}[Fulton--MacPherson Compactification]
\label{def:FM_compactification}
The \emph{Fulton--MacPherson compactification} $\overline{\Conf}_n(\C) := \FM_n(\C)$ is a compact manifold with corners, obtained by adding boundary divisors $D_S$ for each subset $S \subseteq \{1,\ldots,n\}$ with $|S| \geq 2$, parametrizing configurations where points indexed by $S$ approach a common limit (forming a ``cluster'') while remaining separated from points outside $S$.
\end{definition}

\begin{remark}[Local Structure Near Boundaries]
Near a boundary divisor $D_S$, the FM compactification admits local coordinates $(\varepsilon_S, u_S, \theta_S, \ldots)$ where:
\begin{itemize}
\item $\varepsilon_S \geq 0$ is a real radial parameter measuring the ``size'' of cluster $S$ ($\varepsilon_S = 0$ on $D_S$);
\item $u_S \in \C$ is the cluster center;
\item $\theta_S$ are angular variables parametrizing relative positions within the cluster.
\end{itemize}
The boundary divisor $D_S$ is a codimension-$1$ stratum naturally isomorphic to $\FM_{n-|S|+1}(\C) \times \FM_{|S|}(\C)$: the first factor records positions of the $n - |S|$ unclustered points together with the cluster center, and the second factor records the internal configuration of the $|S|$ colliding points at infinitesimal scale. More generally, a deeper boundary stratum corresponding to a partition $\pi = \{S_1, \ldots, S_r\}$ of $\{1,\ldots,n\}$ into $r$ clusters (with $|S_j| \geq 2$ for at least two parts, or nested clusters within a single part) takes the form $\FM_r(\C) \times \prod_{j=1}^r \FM_{|S_j|}(\C)$.
\end{remark}

\subsubsection{Logarithmic Forms and Residues}

The chiral OPE at a collision $z_i \to z_j$ extracts a pole, and a
pole is the residue of a logarithmic $1$-form along the
corresponding boundary divisor of $\FM_n(\C)$. The differential
forms carrying exactly this data are the logarithmic forms in the
Arnold--Orlik--Solomon sense: smooth on the open configuration
stratum, with at most $d\log\varepsilon_S$-type singularities along
each collision divisor $D_S$.

\begin{definition}[Logarithmic Forms]
Let $\Omega^\bullet_{\log}(\FM_n(\C))$ denote the complex of differential forms on $\FM_n(\C)$ with at most \emph{logarithmic singularities} along boundary divisors. These are forms that locally near $D_S$ have the form
\[
\omega = \alpha + \frac{d\varepsilon_S}{\varepsilon_S} \wedge \beta,
\]
where $\alpha, \beta$ extend smoothly to $\varepsilon_S = 0$.
\end{definition}


\begin{definition}[Residue Operator]
\label{def:residue_operator}
For each boundary divisor $D_S$, define the \emph{residue operator}
\[
\Res_{D_S} : \Omega^\bullet_{\log}(\FM_n(\C)) \longrightarrow \Omega^{\bullet-1}_{\log}(D_S)
\]
by writing $\omega = d\log\varepsilon_S \wedge \beta + \alpha$ with $\alpha, \beta$ smooth at $\varepsilon_S = 0$ and setting
\[
\Res_{D_S}(\omega) := \beta\big|_{\varepsilon_S = 0} = \lim_{\varepsilon_S \to 0}\bigl(\varepsilon_S \cdot \iota_{\partial_{\varepsilon_S}} \omega\bigr)\big|_{D_S}.
\]
This extracts the coefficient of the logarithmic $1$-form $d\log\varepsilon_S = d\varepsilon_S/\varepsilon_S$, not the naive contraction $\iota_{\partial_{\varepsilon_S}}\omega|_{D_S}$ (which diverges for logarithmic forms).
\end{definition}


\subsubsection{Construction of the Chiral Operad}

The configuration-space geometry of $\FM_n(\C)$ and the residue
calculus of logarithmic forms combine into a single operad: chains
on $\FM_n(\C)$ with logarithmic coefficients, tensored with iterated
Laurent series in spectral parameters that record derivatives of
the $\lambda$-bracket. Operadic composition is boundary restriction
followed by the chiral product, read off as an OPE residue.

\begin{construction}[Chiral Operad $\cP^{\mathrm{ch}}$]
\label{def:chiral_operad}
The \emph{chiral operad} $\cP^{\mathrm{ch}}$ is defined as follows:

\textbf{Colors:} A single color, representing a generic vertex algebra or chiral algebra.

\textbf{Operations:} For $n \geq 1$, the $n$th component is
\[
\cP^{\mathrm{ch}}(n) := \Omega^{2n-3}_{\log}(\FM_n(\C)) \hat{\otimes} \C((\lambda_1)) \cdots ((\lambda_{n-1})),
\]
where $\FM_n(\C)$ denotes the Fulton--MacPherson compactification modulo translations \emph{and} dilations (real dimension $2n-3$), in contrast to the convention in Section~\ref{sec:FM_calculus} where only translations are quotiented (giving real dimension $2(n-1)$). The forms are top-degree on this moduli space to allow integration over the fundamental class, and $\C((\lambda_1)) \cdots ((\lambda_{n-1}))$ represents iterated Laurent series in spectral parameters $\lambda_i$ (encoding derivatives via generating functions $e^{\lambda(z-w)}$). The comparison theorem (Theorem~\ref{thm:closed-BD}) accommodates this discrepancy by fiber-integrating over the dilation direction $\R_{>0}$: the top-degree forms on $\FM_k(\C)/\text{transl}$ restrict to degree $2k-3$ forms on $\FM_k^{\mathrm{red}}(\C) = \FM_k(\C)/(\text{transl} \times \R_{>0})$ via integration along the dilation fiber.

\textbf{Operadic composition:} Given $\alpha \in \cP^{\mathrm{ch}}(k)$ and $\beta \in \cP^{\mathrm{ch}}(\ell)$, composition $\alpha \circ_i \beta$ is defined by:
\begin{enumerate}[label=(\roman*)]
\item Restrict $\alpha$ to the boundary stratum where the $i$th input collides with the newly inserted cluster, obtaining $\Res_{D_{\{i, n+1\}}}(\alpha) \in \Omega^{2k-4}_{\log}(D_{\{i, n+1\}})$, where $D_{\{i, n+1\}}$ is the divisor parametrizing the collision of the $i$th point with the inserted cluster (requiring $|S| \geq 2$);
\item Pull back $\beta$ to the collision divisor;
\item Wedge the resulting forms and adjust spectral parameters via $\lambda_{\text{new},j} = \lambda_{\text{old},j} + \partial_z$ (translation at cluster center).
\end{enumerate}

\textbf{Completion:} We complete with respect to the pole-order filtration: forms with poles of order $\leq N$ at each divisor form a filtered subsystem.
\end{construction}

\begin{remark}[Dimension conventions for $\FM_k(\C)$]
\label{rem:FM-dimension-conventions}
See Remark~\ref{rem:FM-convention} for the full statement.
In brief: $\FM_k(\C)$ modulo translations has $\dim_\R = 2(k-1)$;
the reduced space $\FM_k^{\mathrm{red}}(\C) = \FM_k(\C)/\R_{>0}$
has $\dim_\R = 2k-3$ (the operadic convention).
The chiral operad (Construction~\ref{def:chiral_operad}) uses
the reduced convention; the passage between conventions is fiber
integration along the dilation direction~$\R_{>0}$.
\end{remark}

\begin{remark}[Relationship to Vertex Algebras]
The chiral operad $\cP^{\mathrm{ch}}$ encodes the operadic structure underlying vertex algebras, including configuration space geometry. The completion and spectral parameters implement the formal calculus of $\lambda$-brackets.
\end{remark}

\subsection{The Two-Colored Operad: Ambient Conventions}
\label{sec:SC-operad-foundations}

The chiral operad $\cP^{\mathrm{ch}}$ of
\S\ref{sec:chiral_operad} governs holomorphic collisions alone. To
couple it to an ordered topological direction we must specify the
target dg-category in which our operads act, and the
topological-to-algebraic functor that converts spaces of
configurations into chain-level operations. Both are standard and
are fixed once for the remainder of the chapter.

We work in a presentable, stable, symmetric monoidal dg-category $(\mathcal C,\otimes,\mathbf 1)$ (e.g.\ $\mathrm{Ch}_k$ for a field $k$ of characteristic zero), with $\otimes$ exact and preserving colimits in each variable.
Topological operads $O$ are replaced by dg-operads $C_\ast(O)$ via singular chains (or de Rham forms if convenient); we silently identify these.

\begin{notation}
\label{not:fm-e1}
Let $\FM_k(\C)$ denote the Fulton--MacPherson compactification of $\Conf_k(\C)$; let $E_1(m)$ be the little-intervals operad. We set $C_\ast(\FM_k(\C))$ and $C_\ast(E_1(m))$ for their chains.
\end{notation}

\subsection{The Two-Coloured Dioperad $\mathsf{SC}^{\mathrm{ch,top}}$}
\label{subsec:sc-def}

$\SCchtop$ is properly a \emph{coloured dioperad} in the sense of
Gan~\cite{Gan03}, not a standard (uncoloured or symmetric-coloured)
operad. The directional colour-restriction is essential: an HT
observable lives on a rectangle $D \times I \subset \C \times \R$,
holomorphic in the $\C$-factor and locally constant in the
$\R$-factor, and the physical coupling is one-way: collision data
flows from the holomorphic bulk ($\mathsf{ch}$) into the topological
boundary ($\mathsf{top}$) but never back. Consequently
$\Sigma_n$-equivariance fails across the colour boundary: a
mixed-colour input tuple $(\mathsf{ch},\ldots,\mathsf{ch},
\mathsf{top},\ldots,\mathsf{top})$ cannot be symmetrised without
breaking the directional restriction, because the permutations that
interchange a $\mathsf{ch}$ input with a $\mathsf{top}$ input produce
operations whose output colour is not preserved. Equivariance holds
separately within each colour class but not between them. This is
exactly the defining content of a dioperad: directed edges in the
underlying composition graph admit colour-restricted (rather than
$\Sigma$-symmetrised) composition. Dunn additivity, which is an
operadic additivity statement for uncoloured little-disk operads,
therefore does \emph{not} apply to $\SCchtop$; the tensor-product
formulas $E_p \otimes E_q \simeq E_{p+q}$ are statements about
symmetric operads and are not available for coloured dioperads with
directional colour-restriction.

The object encoding exactly the holomorphic--topological HT observable
content is built componentwise from $\FM_k(\C)$ and $E_1(m)$, with the
no-open-to-closed rule imposed by setting the would-be ch-producing
mixed operations to the empty space.

\begin{definition}[Colours and spaces of dioperadic operations]
\label{def:SC}
Let $\mathbf{Col}=\{\mathsf{ch},\mathsf{top}\}$. The
\emph{holomorphic--topological Swiss-cheese coloured dioperad} (in the
sense of \cite{Gan03}) $\mathsf{SC}^{\mathrm{ch,top}}$ has spaces of
operations
\begin{align*}
\mathsf{SC}^{\mathrm{ch,top}}\bigl((\underbrace{\mathsf{ch},\dots,\mathsf{ch}}_{k});\mathsf{ch}\bigr) &:= \FM_k(\C),\\
\mathsf{SC}^{\mathrm{ch,top}}\bigl((\underbrace{\mathsf{ch},\dots,\mathsf{ch}}_{k},\underbrace{\mathsf{top},\dots,\mathsf{top}}_{m});\mathsf{top}\bigr) &:= \FM_k(\C)\times E_1(m),\\
\mathsf{SC}^{\mathrm{ch,top}}(\ldots;\mathsf{ch}) &= \varnothing \ \text{ if any input is $\mathsf{top}$}.
\end{align*}
Composition is induced componentwise by the standard insertion maps for $\FM_\bullet(\C)$ and $E_1$, with symmetric-group actions on labeled inputs. Units are the degree–$1$ points in $\FM_1(\C)$ and $E_1(1)$.
\end{definition}

\begin{proposition}[\ClaimStatusProvedHere]
\label{prop:SC-operad}
$\mathsf{SC}^{\mathrm{ch,top}}$ is a well-defined strict
two-coloured topological dioperad in the sense of
Gan~\cite{Gan03}. Associativity, unitality, and within-colour
symmetric-group actions hold strictly; inter-colour equivariance is
not part of the dioperad axiom data (and would conflict with the
directional colour restriction).
\end{proposition}

\begin{proof}
We verify the dioperad axioms directly from the component structure.
\emph{Associativity:} composition in $\SCchtop$ is defined componentwise by FM insertion maps on the $\FM_k(\C)$ factor and interval composition on the $E_1(m)$ factor; since these factors are independent, associativity of the composite follows from associativity of each factor separately.
\emph{Unitality:} the unit in degree~$1$ is the pair $(\mathrm{pt},\mathrm{pt}) \in \FM_1(\C) \times E_1(1)$; both components are the identity element of their respective operads, so the unit axiom holds.
\emph{Within-colour $\Sigma$-equivariance:} the symmetric group $\Sigma_k$ acts on the $\FM_k(\C)$ factor by permuting the labels of input points (with trivial action on $E_1(m)$, whose inputs carry a fixed linear ordering); the equivariance of FM operadic insertion under $\Sigma_k$ implies the within-closed-colour equivariance for $\SCchtop$. No equivariance is postulated across the colour boundary, consistent with the dioperad axiomatic.
\emph{No-open-to-closed stability:} the operation space $\SCchtop(\ldots;\mathsf{ch})$ is empty whenever any input is $\mathsf{top}$; composing any operation into an empty space yields an empty space, so this constraint is preserved under all compositions.
\end{proof}

\subsection{Algebras over $\mathsf{SC}^{\mathrm{ch,top}}$ and its $W$-resolution}
\label{subsec:WSC-alg}
\begin{definition}[Strict and homotopy algebras]
\label{def:SC-alg}
A \emph{strict} $\mathsf{SC}^{\mathrm{ch,top}}$-algebra in $\mathcal C$ is a pair $(A_{\mathsf{ch}},A_{\mathsf{top}})$ with structure maps
\[
C_\ast(\FM_k(\C))\otimes A_{\mathsf{ch}}^{\otimes k}\to A_{\mathsf{ch}},\qquad
C_\ast(\FM_k(\C)\times E_1(m))\otimes A_{\mathsf{ch}}^{\otimes k}\otimes A_{\mathsf{top}}^{\otimes m}\to A_{\mathsf{top}}
\]
satisfying operadic identities. The \emph{homotopy} notion is defined via the Boardman--Vogt resolution $W(\mathsf{SC}^{\mathrm{ch,top}})$: a \emph{$W(\mathsf{SC}^{\mathrm{ch,top}})$–algebra} is an algebra over $C_\ast(W(\mathsf{SC}^{\mathrm{ch,top}}))$, encoding all higher coherences.
\end{definition}

\begin{remark}[Homotopy coherence]
The $W$–resolution equips trees indexing compositional patterns with edge-length parameters, turning strict equalities into explicit higher homotopies. This is the standard device to encode $A_\infty$/$E_1$ coherence in the operadic setting.
\end{remark}

\subsection{Closed-Color Comparison with the BD Chiral Operad}
\label{subsec:BD-compare}

Beilinson--Drinfeld construct the chiral operad $\mathsf P_{\mathrm{ch}}$ abstractly, via $\ast$-operations on $D$-modules on powers of the curve~\cite[\S3.4]{BD04}. Our construction builds the closed color of $\SCchtop$ concretely, from logarithmic forms on the FM compactification. The two must agree up to quasi-isomorphism; the comparison proceeds through the Ran space.

For the statement, let $\mathscr L_{\log}$ denote the \emph{logarithmic local system} on $\FM_n(\C)$: the $\Z$-graded constructible system whose stalk over a boundary divisor $D_S$ records the $d\log\varepsilon_S$-pole data encoded by the residue operator (Definition~\ref{def:residue_operator}), with monodromy recording the chiral OPE residues.

\begin{theorem}[Closed color comparison; \ClaimStatusProvedHere]
\label{thm:closed-BD}
There exists a zigzag of quasi-isomorphisms of operads between the BD chiral operad $\mathsf P_{\mathrm{ch}}$ and the dg-operad $C_\ast(\FM_\bullet(\C);\mathscr L_{\log})$ of chains on $\FM_\bullet(\C)$ with logarithmic local system $\mathscr L_{\log}$. Hence, restricting $W(\mathsf{SC}^{\mathrm{ch,top}})$ to the closed color recovers $\mathsf P_{\mathrm{ch}}$ up to weak equivalence.
\end{theorem}

\begin{proof}
The argument proceeds in three steps, constructing a zigzag through an intermediate category of factorization algebras on the Ran space.

\medskip\noindent
\textbf{Step 1: The Ran space and factorization algebras.}
Consider the Ran space
\[
\Ran(\C) \;=\; \operatorname{colim}_{S \in \mathrm{Fin}^{\mathrm{surj}}} \C^S,
\]
the space of finite non-empty subsets of $\C$, topologized so that points can collide but not separate. A \emph{factorization algebra} on $\C$ is a cosheaf $\mathcal{F}$ on $\Ran(\C)$ equipped with factorization isomorphisms
\[
\mathcal{F}(S_1 \sqcup S_2) \;\xrightarrow{\;\sim\;}\; \mathcal{F}(S_1) \otimes \mathcal{F}(S_2)
\]
for disjoint subsets, satisfying associativity and locality constraints. Denote by $\Fact(\C)$ the $(\infty,1)$-category of factorization algebras on $\C$.

\medskip\noindent
\textbf{Step 2: The zigzag.}
We construct a zigzag of quasi-isomorphisms of operads
\[
C_\ast(\FM_\bullet(\C);\mathscr{L}_{\log}) \;\xleftarrow{\;\;\alpha\;\;}\; \Fact(\C)^{\otimes} \;\xrightarrow{\;\;\beta\;\;}\; \mathsf{P}_{\mathrm{ch}}.
\]

\emph{Arrow $\alpha$ (left leg):} The FM compactification $\FM_n(\C)$ provides a model for the configuration space strata of $\Ran(\C)$. The logarithmic local system $\mathscr{L}_{\log}$ on $\FM_n(\C)$ encodes the meromorphic OPE poles: near the boundary divisor $D_S$ where points indexed by $S$ collide, the logarithmic $1$-form $d\log\varepsilon_S$ captures the pole $z_i \to z_j$. Integration of factorization algebra structure maps over the fundamental classes $[\FM_n(\C)]$ with coefficients in $\mathscr{L}_{\log}$ yields operadic compositions in $C_\ast(\FM_\bullet(\C);\mathscr{L}_{\log})$. This map is a quasi-isomorphism for two reasons: first, $\FM_n(\C)$ is a smooth compactification of $\Conf_n(\C)/\mathrm{transl}$ whose inclusion $\Conf_n(\C)/\mathrm{transl} \hookrightarrow \FM_n(\C)$ is a homotopy equivalence (the added boundary strata are contractible normal bundles), so $\FM_n(\C)$ is a homotopy-theoretic model for the $n$-point stratum of $\Ran(\C)$; second, the logarithmic local systems on $\FM_n(\C)$ encode precisely the singular behaviour permitted by the factorization locality axiom, namely the meromorphic OPE poles at collision diagonals, and the de~Rham cohomology of $\FM_n(\C)$ with coefficients in $\mathscr{L}_{\log}$ computes the factorization homology of the corresponding stratum. The construction follows \cite[{\S}4.5.3]{BD04}, which identifies chiral operations on the Ran space with the combinatorics of FM boundary strata.

\emph{Arrow $\beta$ (right leg):} Beilinson--Drinfeld's main theorem \cite[{\S}4.5.3]{BD04} establishes that the pseudo-tensor category of chiral algebras (formulated via $D$-modules on $\Ran(\C)$) is equivalent to the category of factorization algebras on $\C$. Concretely, the BD chiral operad $\mathsf{P}_{\mathrm{ch}}$ governs chiral algebras via the $\ast$-operations $\mu^{(n)}\colon j_\ast j^\ast(\mathcal{A}^{\boxtimes n}) \to \Delta_\ast \mathcal{A}$ on the $n$th power of $\C$. These $\ast$-operations correspond, under the BD equivalence, to the factorization structure maps evaluated at colliding configurations. The derived/homotopical refinement is due to Francis--Gaitsgory \cite[Theorem~5.1]{FG12}, who promote this to a quasi-isomorphism $\beta$ of operads in the $(\infty,1)$-categorical setting.

\medskip\noindent
\textbf{Step 3: Conclusion.}
Both arrows $\alpha$ and $\beta$ are quasi-isomorphisms of operads. The claimed weak equivalence is the composition $\beta \circ \alpha^{-1}$ (well-defined in the homotopy category of operads, since $\alpha$ is an equivalence), giving
\[
C_\ast(\FM_\bullet(\C);\mathscr{L}_{\log}) \;\simeq\; \mathsf{P}_{\mathrm{ch}}.
\]
Since restricting $\SCchtop$ to the closed color $\mathsf{ch}$ yields precisely the operad with $n$th component $\FM_n(\C)$ (Definition~\ref{def:SC}), and the Boardman--Vogt resolution $W(-)$ preserves weak equivalences of well-pointed $\Sigma$-cofibrant operads, restricting $W(\SCchtop)$ to the closed color recovers $\mathsf{P}_{\mathrm{ch}}$ up to weak equivalence.
\end{proof}

\subsection{Recognition: local HT shadows and $W(\SCchtop)$-algebras}
\label{subsec:SC-vs-prefactorization}

The closed-color comparison handles holomorphic data alone. A
3d HT observable has two colors coupled: holomorphic along
$\C$, topological along $\R$, with OPE residues flowing from bulk
to boundary. To get a local-to-global principle that sees both
colors we work directly with observables on rectangles
$D \times I \subset \C \times \R$. The output of this local analysis
is the formal collision shadow of the factorization algebra: a
$W(\SCchtop)$-algebra. It is not the whole Ran-space or D-module
factorization datum.

\begin{definition}[HT prefactorization]
\label{def:prefact}
A \emph{holomorphic--topological prefactorization algebra} on $\C\times\R$ is a functor $\Obs$ from disjoint unions of open rectangles $D\times I$ to $\mathcal C$, with structure maps $\bigotimes\Obs(U_i)\to \Obs(\bigsqcup U_i)$, satisfying: holomorphic factorization along~$\C$, local constancy along~$\R$, and excision/gluing coherent with shape composition.
\end{definition}

The assignment $U \mapsto \Obs(U)$ is defined only on rectangles.
To extract the product-formal local operadic shadow one needs a
descent principle on product covers. Weiss cosheaf descent is the
correct principle for factorization algebras, but the argument below
uses only product covers and does not claim descent for arbitrary
rectangle Weiss covers.

\begin{definition}[Weiss covers and cosheaf descent]
\label{def:weiss-cover}
Let $M$ be a smooth manifold and $\mathsf{Disk}(M)$ the poset of
open subsets diffeomorphic to disjoint unions of disks (or
rectangles, for product manifolds). A \emph{Weiss cover} of
$U \in \mathsf{Disk}(M)$ is a collection $\{U_\alpha\}_{\alpha \in A}$
of opens in $\mathsf{Disk}(M)$ contained in~$U$ such that for every
finite set of points $S \subset U$ there exists~$\alpha$ with
$S \subset U_\alpha$. A prefactorization algebra $\Obs$ on~$M$
satisfies \emph{Weiss cosheaf descent} if for every
$U \in \mathsf{Disk}(M)$ and every Weiss cover
$\{U_\alpha\}$ of~$U$, the natural map
\[
 \operatorname{hocolim}_{\check{C}(\{U_\alpha\})} \Obs
 \;\xrightarrow{\;\sim\;}\;
 \Obs(U)
\]
is a quasi-isomorphism, where
$\check{C}(\{U_\alpha\})$ denotes the \v{C}ech nerve of the cover
(formed using pairwise intersections in $\mathsf{Disk}(M)$).
\end{definition}

\begin{lemma}[Product-cover Weiss descent for $\C \times \R$; \ClaimStatusProvedHere]
\label{lem:product-weiss-descent}
Let $\Obs$ be an HT prefactorization algebra on~$\C \times \R$
(Definition~\textup{\ref{def:prefact}}) that is holomorphic along~$\C$ and
locally constant along~$\R$. Let
$\{D_\beta\}_{\beta\in B}$ be a Weiss cover of a disk $D\subset\C$
and let $\{I_\gamma\}_{\gamma\in G}$ be a Weiss cover of an interval
$I\subset\R$. Then the product cover
$\{D_\beta\times I_\gamma\}_{(\beta,\gamma)\in B\times G}$ computes
$\Obs(D\times I)$ by Weiss cosheaf descent.
\end{lemma}

\begin{proof}
The proof applies only to product covers. It does not assert that an
arbitrary Weiss cover by rectangles is cofinally refined by a product
cover.

\medskip\noindent
\textbf{Step~1: Holomorphic descent along $\C$.}
Costello--Gwilliam~\cite[Chapter~5, \S5.4]{CG17} prove that a
holomorphic prefactorization algebra on~$\C$ with meromorphic OPE
singularities satisfies Weiss cosheaf descent: the
holomorphic locality condition (factorization along disjoint disks,
with controlled OPE poles along collision divisors)
implies that $D \mapsto \Obs(D \times I)$ satisfies descent with
respect to any Weiss cover of~$D$ by sub-disks.
Concretely, the argument uses the fact that meromorphic forms on
$\FM_k(\C)$ extend to the compactification and that the
factorization isomorphisms
$\Obs(D_1 \sqcup D_2) \xrightarrow{\sim} \Obs(D_1) \otimes
\Obs(D_2)$, together with the OPE residue maps on boundary strata,
provide the gluing data for the \v{C}ech nerve. This identifies the
holomorphic factor with an algebra over
$C_\ast(\FM_\bullet(\C); \mathscr{L}_{\log})$.

\medskip\noindent
\textbf{Step~2: Topological descent along $\R$.}
The Ayala--Francis recognition theorem~\cite[Theorem~3.24]{AF15} applies to the
$\R$-direction. Their theorem states: for a manifold~$M$ of
dimension~$n$, the functor from $\mathsf{Disk}_n$-algebras (i.e.\
$E_n$-algebras) to locally constant factorization algebras on~$M$ is
an equivalence when $M$ is framed. For $M = \R$ ($n = 1$, canonically
framed), this says: locally constant prefactorization algebras
on~$\R$ satisfying Weiss descent are equivalent to $E_1$-algebras.

Since $\Obs$ is locally constant along~$\R$ by hypothesis,
$I \mapsto \Obs(D \times I)$ is locally constant: for any inclusion
$I' \hookrightarrow I$ of intervals, the restriction map
$\Obs(D \times I') \to \Obs(D \times I)$ is a quasi-isomorphism.
Locally constant prefactorization algebras on~$\R$ automatically
satisfy Weiss descent (a Weiss cover of an interval~$I$ by
sub-intervals $\{I_\alpha\}$ has a contractible \v{C}ech nerve, and
local constancy implies that all restriction maps in the nerve are
quasi-isomorphisms, so the hocolimit computes $\Obs(D \times I)$).
This identifies the topological factor with an $E_1$-algebra.

\medskip\noindent
\textbf{Step~3: Product-cover descent.}
For the product cover $\{D_\beta\times I_\gamma\}$, the
\v{C}ech nerve is the product of the two factor nerves. Hence the
hocolimit over the product cover factors as
\[
 \operatorname{hocolim}_{\check{C}(\{D_\beta \times I_\gamma\})} \Obs
 \;\simeq\;
 \operatorname{hocolim}_{\check{C}(\{D_\beta\})}
 \Bigl(D_\beta \mapsto
 \operatorname{hocolim}_{\check{C}(\{I_\gamma\})} \Obs(D_\beta \times -)
 \Bigr).
\]
The inner hocolimit (over~$I$) is a quasi-isomorphism
$\xrightarrow{\sim} \Obs(D_\beta \times I)$ by Step~2
(topological descent). The outer hocolimit (over~$D$) is then a
quasi-isomorphism
$\xrightarrow{\sim} \Obs(D \times I)$ by Step~1
(holomorphic descent).
Composing, the full Weiss descent map is a quasi-isomorphism.
\end{proof}

\begin{remark}[Role of the product structure]
\label{rem:product-weiss}
The proof of Lemma~\ref{lem:product-weiss-descent} uses the product
structure $\C \times \R$ in an essential way: the factorization of
the iterated hocolimit argument relies on an actual product cover,
not on an arbitrary rectangle Weiss cover. For arbitrary rectangle
covers, or for a general $3$-manifold~$M$ that is not a product, the
Weiss descent comparison requires a separate cofinality or
factorization-envelope argument (e.g.\ as in~\cite{CG17}).
The product structure is the geometric reflection of the physical
fact that holomorphic and topological directions decouple at the level
of the underlying manifold (Remark~\ref{rem:product-forced-by-locality}
in Section~\ref{sec:foundations}).
\end{remark}

\begin{theorem}[Local Swiss-cheese shadow; \ClaimStatusConditional]
\label{thm:recognition-SC}
Let $\Obs$ be an HT prefactorization algebra on the rectangle basis
of $\C\times\R$ satisfying product Weiss descent
\textup{(}Lemma~\ref{lem:product-weiss-descent}\textup{)}. Formal
completion at collisions in the holomorphic direction and at ordered
interval collisions in the topological direction assigns to $\Obs$ a
canonical $C_\ast(W(\mathsf{SC}^{\mathrm{ch,top}}))$-algebra
$(A_{\mathsf{ch}},A_{\mathsf{top}})$.

Conversely, a $C_\ast(W(\mathsf{SC}^{\mathrm{ch,top}}))$-algebra
has a canonical rectangle factorization envelope
$\operatorname{Env}_{\mathrm{rect}}(A_{\mathsf{ch}},A_{\mathsf{top}})$ generated
by its two colors. These two constructions form an adjunction. The
unit and counit are equivalences precisely on the product-formal
rectangle prefactorization algebras generated by their two local
stalk colors.
\end{theorem}

\begin{proof}
The proof identifies the local operadic shadow and records the
essential-image restriction explicitly.

\medskip\noindent
\textbf{Step 1: The local shadow of rectangle decompositions.}
An HT prefactorization algebra $\Obs$ assigns a cochain complex to
each open rectangle $D \times I \subset \C \times \R$, with structure
maps for disjoint inclusions. A configuration of holomorphic
sub-disks and ordered topological sub-intervals determines a
rectangle decomposition:
\begin{itemize}
\item Holomorphic disks $D_1,\ldots,D_k\hookrightarrow D$ determine
 a point of $\FM_k(\C)$ after passing to the formal completion at the
 collision stratum.
\item Ordered intervals $I_1,\ldots,I_m\hookrightarrow I$ determine
 a point of $E_1(m)$, recording only the order and nesting data.
\end{itemize}
Thus the formal local operation space is
$\FM_k(\C)\times E_1(m)$, the mixed operation space of
$\SCchtop$. The no-open-to-closed rule is inherited from the HT
directionality: topological boundary inputs do not create a
holomorphic closed output.

\medskip\noindent
\textbf{Step 2: $W$-coherences and homotopy-coherent associativity.}
The Boardman--Vogt resolution $W(\SCchtop)$ decorates each internal
edge $e$ of a composition tree with a parameter $s_e\in[0,1]$. These
parameters encode the homotopies between different orders of
performing nested rectangle decompositions:
\begin{itemize}
\item $s_e=0$: the two operations sharing edge $e$ are composed in
 one step.
\item $s_e=1$: the inner and outer decompositions are performed
 sequentially.
\end{itemize}
The cubical parameter space $[0,1]^{|E_{\mathrm{int}}|}$ supplies the
higher homotopies for iterated inclusions. This is exactly the
standard Boardman--Vogt construction~\cite{BV73}; applying chains
gives coherent $C_\ast(W(\SCchtop))$ operations.

\medskip\noindent
\textbf{Step 3: From a prefactorization algebra to its operadic shadow.}
Choose sufficiently small base rectangles around a point
$(z_0,t_0)$. Local constancy in the $\R$-direction and holomorphic
factorization in the $\C$-direction identify all smaller choices up
to canonical quasi-isomorphism; equivalently one may take the
corresponding germs. The closed color $A_{\mathsf{ch}}$ is the germ
seen by holomorphic collision operations, and the open color
$A_{\mathsf{top}}$ is the germ seen by ordered interval operations.

For each configuration in $\FM_k(\C)\times E_1(m)$, the
prefactorization structure map of $\Obs$ on the corresponding
rectangle decomposition gives a multilinear operation on these two
germs. Passing to chains and inserting the $W$-edge parameters from
Step~2 produces the required operation
\[
 C_\ast(W(\SCchtop))((\mathsf{ch}^k,\mathsf{top}^m);\mathsf{out})
 \otimes A_{\mathsf{ch}}^{\otimes k}
 \otimes A_{\mathsf{top}}^{\otimes m}
 \longrightarrow A_{\mathsf{out}} .
\]
The operadic identities are the functoriality and associativity of
prefactorization structure maps, restricted to boundary strata of
$\FM_k(\C)\times E_1(m)$. Product Weiss descent ensures that replacing
a rectangle by a Weiss refinement does not change the resulting germ.

\medskip\noindent
\textbf{Step 4: The rectangle factorization envelope.}
Conversely, let $(A_{\mathsf{ch}},A_{\mathsf{top}})$ be a
$C_\ast(W(\SCchtop))$-algebra. Define
$\operatorname{Env}_{\mathrm{rect}}(A_{\mathsf{ch}},A_{\mathsf{top}})$ on a
rectangle by the left Kan extension from the full subcategory of
standard finite unions of formal disks and formal intervals. On a
standard configuration its value is the ordered tensor product of the
corresponding color labels, modulo the relations imposed by the
$W(\SCchtop)$-operations. The factorization maps are induced by
operadic insertion. Since $W(\SCchtop)$ is cofibrant, these maps carry
the needed coherent homotopies.

This construction is the usual factorization envelope on the rectangle
basis. It is left adjoint to taking the local operadic shadow:
a map from the envelope to a rectangle prefactorization algebra is the
same datum as compatible images of the two colors satisfying the
$W(\SCchtop)$ operation relations.

\medskip\noindent
\textbf{Step 5: Essential image.}
The previous two steps do not prove that every HT prefactorization
algebra is recovered from its local operadic shadow. They prove
recovery only for prefactorization algebras generated from the two
stalk colors by rectangle factorization and product Weiss descent.
Global Ran-space data, D-module monodromy, and higher-genus curvature
belong to the factorization layer and are invisible to
$W(\SCchtop)$. This is the local/global distinction used in
Chapter~\ref{ch:factorization-swiss-cheese}.
\qedhere
\end{proof}

\begin{corollary}[Clean replacement of the semidirect slogan;
\ClaimStatusConditional]
\label{cor:clean-replacement}
The two--colored operad $W(\mathsf{SC}^{\mathrm{ch,top}})$ replaces the ambiguous $W(\mathsf P_{\mathrm{ch}}\rtimes E_1)$ and precisely governs the local formal datum ``chiral in~$z$ + $E_1$ along~$t$ + all rectangle-level compatibilities.'' The product decomposition $\FM_k(\C) \times E_1(m)$ of the operation spaces, the no-open-to-closed rule, and the Boardman--Vogt coherences determine that local operad uniquely. Global Ran-space and monodromy data remain part of the factorization layer, not of this semidirect replacement.
\end{corollary}

\begin{remark}[Ribbon/trace enhancement for literal 't~Hooft sums]%
\label{rem:ribbon-trace-enhancement}%
\index{ribbon structure!necessity for t Hooft@necessity for 't~Hooft}%
\index{t Hooft expansion@'t~Hooft expansion!ribbon enhancement}%
The tree-level two-colored operad
$W(\mathsf{SC}^{\mathrm{ch,top}})$ is the correct
genus-$0$ Swiss-cheese object. For literal large-$N$
't~Hooft counting, two additional structures are needed
\textup{(}Vol~I,
Cyclicity-Ribbon Theorem and
Bare-Graph Lemma\textup{)}:
\begin{enumerate}[label=\textup{(\roman*)}]
\item a \emph{cyclic trace} $\tau$ on the open
 $E_1$-color, turning linear vertices into cyclic
 tensors and enabling canonical ribbon thickening;
\item a \emph{matrix realization}
 $O_N = \mathrm{Mat}_N \otimes O_0$, so that each
 face of the thickened surface contributes a factor
 of~$N$ via index-loop summation.
\end{enumerate}
The resulting object is the \emph{ribbonized modular
Swiss-cheese theory}
$\mathrm{RSC}_{\mathrm{tr}}^{\mathrm{mod}}$
\textup{(}Vol~I,
Definition of ribbonized Swiss-cheese\textup{)},
whose graph amplitudes carry the exact weight
$N^{\chi(\Sigma_\Gamma)}$ \textup{(}Vol~I,
Exact $N^{\chi}$-Weighting Theorem\textup{)}.
Without this enhancement, bare stable-graph sums give
genus-organized expansions but not canonical
$N$-weighted 't~Hooft sums.
\end{remark}

\begin{remark}[HT locality and the universal-holography climax]
\label{rem:locality-holography-anchor}%
\index{HT locality!universal-holography anchor}%
The HT prefactorization recognition theorem
(Theorem~\ref{thm:recognition-SC}) supplies the local
$W(\SCchtop)$ shadow used by the universal-holography master theorem
(Theorem~\ref{thm:universal-holography-master}). The product Weiss
descent (Lemma~\ref{lem:product-weiss-descent}) decouples
holomorphic locality on $\C$ from $\Eone$-topological local
constancy on $\R$ at the prefactorization-data level; the
recognition theorem extracts the formal collision object via the
Boardman--Vogt resolution of $\SCchtop$, producing a
$C_*(W(\SCchtop))$-algebra structure on the local boundary
observables. It does not recover global factorization data from the
local operad.
For an algebra $\cA$ with conformal vector at non-critical level,
the resulting $\SCchtop$-algebra topologizes via Sugawara
$T(z) = [Q_{\mathrm{tot}}, G(z)]$ (making $\C$-translations
$Q$-exact) to an $E_3$-topological factorization algebra on
$\C \times \R$; this is the local content of the climax bulk
identification
$\Obs^{\mathrm{bulk}}(T_\cA) \simeq \cZ^{\mathrm{der}}_{\mathrm{ch}}(\cA)$.
At critical level $k = -h^\vee$ the Sugawara construction
degenerates, topologization fails, and the $\SCchtop$ algebra
remains in the generic regime: the Feigin--Frenkel stratum of the
infinite-fingerprint classification, $\kappa = 0$.
The directionality $\SCchtop \neq E_3$ is enforced throughout: the
present recognition theorem produces $\SCchtop$ structure, never
$E_3$ structure directly; the topologization step is a separate
theorem with explicit hypotheses (conformal vector, non-critical
level, chain-level $Q$-exactness witness), proved for affine
Kac--Moody at non-critical level in
Theorem~\ref{thm:E3-topological-km}, conjectural for general
chiral algebras with conformal vector, and chain-level for class
M in the weight-completed category via
Theorem~\ref{thm:programme-climax}.
\end{remark}
